\chapter{为什么共享状态需要同步与通信对象}

\begin{statebox}
\textbf{项目里程碑 v0.10。} 内核已经有可抢占 task、受控入口和中断驱动的设备事件。
只要任务与 IRQ 都能修改同一状态，“各自执行正确”就不再保证组合结果正确；只禁止并发
又会让等待设备或管道数据的任务占住处理器。
\end{statebox}

\inputchapterbody{chapters/ch07-synchronization-ipc-locks}

\section{把同步规则落实为状态机}
\inputdetail{chapters/ch07-stage-two-sync-state-machines}
\inputdetail{chapters/ch07-concurrency-ipc-signals}
\inputdetail{chapters/ch08-kernel-locking-deadlocks}

\begin{problemframe}
\textbf{尚未解决的问题。} 执行与并发边界已经形成，但程序仍只能直接面向驱动协议，
不同设备的请求粒度、队列与持久化命令也无法由进程稳定持有。下一章沿 v0.11 先建立公共
块请求层，再建立 fd、打开文件对象与 VFS 分发。
\end{problemframe}
